\documentclass[11pt,a4paper]{article}
\usepackage[margin=1in]{geometry}
\usepackage{mathptmx}
\usepackage{amsmath}
\usepackage{amssymb}
\usepackage{amsthm}
\usepackage{braket}
\usepackage{tikz}
\usepackage{quantikz}
\usepackage{enumitem}
\usepackage{hyperref}
\usepackage{xcolor}

\definecolor{circuitblue}{RGB}{40,80,160}
\hypersetup{colorlinks=true,linkcolor=circuitblue,urlcolor=circuitblue,citecolor=circuitblue}

\theoremstyle{definition}
\newtheorem{definition}{Definition}[section]
\newtheorem{example}{Example}[section]

\theoremstyle{remark}
\newtheorem*{intuition}{Intuition}
\newtheorem*{keypoint}{Key Point}

\newcounter{exercise}[section]
\newenvironment{exercise}{\refstepcounter{exercise}\par\medskip\noindent\textbf{Exercise \theexercise.}\rmfamily}{\medskip}

\pagestyle{plain}
\setlength{\parindent}{0pt}
\setlength{\parskip}{6pt}

\begin{document}

{\LARGE \textbf{Lab 5: Quantum Classifier}}\\[8pt]
{\large QML}\\[12pt]
\hrule

\vspace{16pt}

\section{introduction}

this is it --- the capstone lab. we r building actual quantum machine learning classifiers and comparing them head-to-head with the classical baselines from lab 4. u'll implement two approaches:

\begin{enumerate}[leftmargin=*]
  \item \textbf{variational quantum classifier (VQC)} --- a parameterized quantum circuit trained with gradient descent, similar to a neural network
  \item \textbf{quantum support vector classifier (QSVC)} --- uses a quantum kernel computed on a quantum computer, plugged into a classical SVM
\end{enumerate}

both approaches use quantum feature maps to encode classical data into quantum states. the question is: does the quantum feature map provide any advantage over classical feature maps?

by the end of this lab u should be able to:
\begin{itemize}[leftmargin=*]
  \item build quantum feature maps (data encoding circuits)
  \item implement a VQC from scratch using qiskit
  \item compute quantum kernel matrices
  \item train a QSVC using the precomputed kernel workflow from lab 4
  \item rigorously compare quantum vs classical classifiers
\end{itemize}

\begin{keypoint}
we r using the same iris dataset (binary, 2 features) as lab 4. same preprocessing, same train/test split. this way the comparison is fair.
\end{keypoint}

\section{quantum feature maps}

\subsection{encoding classical data}

the first step in any quantum ML pipeline: encode classical data $x \in \mathbb{R}^d$ into a quantum state $\ket{\phi(x)}$. this is done by a parameterized circuit where the parameters r set by the data.

\begin{definition}
a \textbf{quantum feature map} is a unitary $U(x)$ that maps classical data $x$ to a quantum state:
\[
\ket{\phi(x)} = U(x)\ket{0}^{\otimes n}
\]
\end{definition}

the most common feature map is the \textbf{ZZ feature map}. for 2 qubits and 2 features $(x_1, x_2)$:

\begin{center}
\begin{quantikz}
\lstick{$q_0$} & \gate{H} & \gate{R_z(x_1)} & \ctrl{1} & \qw & \ctrl{1} & \qw \\
\lstick{$q_1$} & \gate{H} & \gate{R_z(x_2)} & \targ{} & \gate{R_z(x_1 x_2)} & \targ{} & \qw
\end{quantikz}
\end{center}

the H gates create superposition. the $R_z(x_i)$ gates encode individual features. the CNOT-$R_z$-CNOT block encodes the interaction $x_1 \cdot x_2$. this is important bcs it means the feature map captures correlations between features, not just individual features.

another option is the \textbf{angle encoding} feature map, which is simpler:

\begin{center}
\begin{quantikz}
\lstick{$q_0$} & \gate{R_y(x_1)} & \gate{R_z(x_1)} & \qw \\
\lstick{$q_1$} & \gate{R_y(x_2)} & \gate{R_z(x_2)} & \qw
\end{quantikz}
\end{center}

and u can repeat the feature map multiple times (called ``data re-uploading'') to make it more expressive:

\begin{center}
\begin{quantikz}
\lstick{$q_0$} & \gate{H} & \gate{R_z(x_1)} & \ctrl{1} & \qw & \ctrl{1} & \gate{H} & \gate{R_z(x_1)} & \ctrl{1} & \qw & \ctrl{1} & \qw \\
\lstick{$q_1$} & \gate{H} & \gate{R_z(x_2)} & \targ{} & \gate{R_z(x_1 x_2)} & \targ{} & \gate{H} & \gate{R_z(x_2)} & \targ{} & \gate{R_z(x_1 x_2)} & \targ{} & \qw
\end{quantikz}
\end{center}

\begin{exercise}
\textbf{building feature maps.}
\begin{enumerate}[leftmargin=*]
  \item implement the ZZ feature map for 2 qubits manually:
    \begin{itemize}
      \item create a \texttt{QuantumCircuit(2)}
      \item apply H to both qubits
      \item apply $R_z(x_1)$ to qubit 0 and $R_z(x_2)$ to qubit 1
      \item apply CNOT(0,1), then $R_z(x_1 \cdot x_2)$ on qubit 1, then CNOT(0,1)
    \end{itemize}
  \item also implement the angle encoding feature map
  \item test both with a sample data point, e.g., $x = (0.5, -0.3)$
  \item use the statevector simulator to get $\ket{\phi(x)}$ for each feature map
  \item try a few different data points and see how the output states differ
\end{enumerate}

u can also use qiskit's built-in feature maps: \texttt{ZZFeatureMap} and \texttt{ZFeatureMap} from \texttt{qiskit.circuit.library}. compare ur manual implementation with the built-in one.
\end{exercise}

\section{quantum kernel method (QSVC)}

\subsection{the quantum kernel}

the quantum kernel between two data points $x$ and $x'$ is:
\[
k(x, x') = |\braket{\phi(x)|\phi(x')}|^2
\]

this measures how similar the quantum states r for two different data points. to compute this on a quantum computer, we use the following circuit:

\begin{center}
\begin{quantikz}
\lstick{$\ket{0}$} & \gate[2]{U(x)} & \gate[2]{U^\dagger(x')} & \meter{} \\
\lstick{$\ket{0}$} & \qw & \qw & \meter{}
\end{quantikz}
\end{center}

u apply $U(x)$ to encode the first data point, then $U^\dagger(x')$ to ``undo'' the second data point. then u measure: the probability of getting all zeros is exactly $|\braket{\phi(x)|\phi(x')}|^2$.

\begin{intuition}
if $x = x'$, the circuit does $U(x)U^\dagger(x) = I$, so u always measure $\ket{00\ldots0}$ and the kernel is 1. if $x$ and $x'$ r very different, the states will be orthogonal and the kernel will be close to 0. so the kernel measures similarity in quantum state space.
\end{intuition}

\begin{exercise}
\textbf{computing the quantum kernel.}
\begin{enumerate}[leftmargin=*]
  \item write a function \texttt{quantum\_kernel(x, x\_prime, feature\_map)} that:
    \begin{itemize}
      \item applies the feature map $U(x)$
      \item applies the inverse feature map $U^\dagger(x')$
      \item measures all qubits
      \item returns the probability of the all-zeros outcome
    \end{itemize}
  \item test it: compute $k(x, x)$ for some point $x$. u should get $\approx 1$.
  \item compute $k(x, x')$ for two very different points. should be much less than 1.
  \item now write a function \texttt{quantum\_kernel\_matrix(X1, X2, feature\_map)} that computes the full kernel matrix. this requires $|X1| \times |X2|$ circuit evaluations.
  \item compute the kernel matrix for the training data. visualize it as a heatmap. compare with the RBF kernel heatmap from lab 4.
\end{enumerate}
\end{exercise}

\begin{exercise}
\textbf{training QSVC.}
\begin{enumerate}[leftmargin=*]
  \item compute the quantum kernel matrix for training data: $K_{\text{train}}(i,j) = k(x_i, x_j)$
  \item compute the quantum kernel matrix for test data: $K_{\text{test}}(i,j) = k(x_i^{\text{test}}, x_j^{\text{train}})$
  \item use \texttt{SVC(kernel='precomputed')} from sklearn (exactly like exercise in lab 4)
  \item train: \texttt{.fit(K\_train, y\_train)}
  \item predict: \texttt{.predict(K\_test)}
  \item report test accuracy
  \item compare with ur classical SVM baseline from lab 4
\end{enumerate}

use the ZZ feature map with 1 repetition first. then try 2 repetitions. does more repetitions help?
\end{exercise}

\section{variational quantum classifier (VQC)}

\subsection{the architecture}

a VQC has three parts:
\begin{enumerate}[leftmargin=*]
  \item \textbf{feature map} $U(x)$ --- encodes the input data (same as before)
  \item \textbf{variational ansatz} $V(\theta)$ --- parameterized circuit with trainable parameters
  \item \textbf{measurement} --- measure one qubit and interpret as the class label
\end{enumerate}

the full circuit for a 2-qubit VQC:

\begin{center}
\begin{quantikz}
\lstick{$\ket{0}$} & \gate[2]{U(x)} & \gate{R_y(\theta_1)} & \ctrl{1} & \gate{R_y(\theta_3)} & \meter{} \\
\lstick{$\ket{0}$} & \qw & \gate{R_y(\theta_2)} & \targ{} & \gate{R_y(\theta_4)} & \qw
\end{quantikz}
\end{center}

the feature map encodes the data, the ansatz has trainable parameters $\theta$, and we measure qubit 0. if $P(\ket{0}) > 0.5$, predict class 0; otherwise predict class 1.

the ansatz can have multiple layers. a single layer typically has $R_y$ rotations on each qubit followed by CNOT entangling gates. here's a 2-layer ansatz:

\begin{center}
\begin{quantikz}
\lstick{$q_0$} & \gate{R_y(\theta_1)} & \ctrl{1} & \gate{R_y(\theta_3)} & \ctrl{1} & \qw \\
\lstick{$q_1$} & \gate{R_y(\theta_2)} & \targ{} & \gate{R_y(\theta_4)} & \targ{} & \qw
\end{quantikz}
\end{center}

\subsection{training}

training a VQC is an optimization problem:
\[
\min_\theta \mathcal{L}(\theta) = \min_\theta \frac{1}{m}\sum_{i=1}^{m} \ell\big(f_\theta(x_i),\; y_i\big)
\]

where $f_\theta(x)$ is the VQC's prediction and $\ell$ is a loss function (e.g., cross-entropy or MSE).

the challenge: computing gradients of quantum circuits. we use the \textbf{parameter shift rule}:
\[
\frac{\partial}{\partial \theta_k} f(\theta) = \frac{1}{2}\left[f\left(\theta_k + \frac{\pi}{2}\right) - f\left(\theta_k - \frac{\pi}{2}\right)\right]
\]

this means: to compute the gradient w.r.t.\ one parameter, u run the circuit twice (with the parameter shifted by $\pm\pi/2$). for $p$ parameters, u need $2p$ circuit evaluations per gradient step.

\begin{exercise}
\textbf{building the VQC.}
\begin{enumerate}[leftmargin=*]
  \item build the VQC circuit manually:
    \begin{itemize}
      \item create a parameterized circuit with qiskit \texttt{Parameter} objects
      \item first apply the ZZ feature map (with data parameters)
      \item then apply a variational ansatz: 2 layers of $R_y$ rotations + CNOT entangling
      \item add measurement on qubit 0
    \end{itemize}
  \item write a function \texttt{predict(x, theta)} that:
    \begin{itemize}
      \item binds the data parameters to $x$ and the variational parameters to $\theta$
      \item runs the circuit on the qasm simulator
      \item returns $P(\ket{0})$ as the prediction
    \end{itemize}
  \item test with random $\theta$ values. the predictions should be between 0 and 1.
\end{enumerate}

the circuit should look like this (feature map + 1-layer ansatz + measurement):

\begin{center}
\begin{quantikz}
\lstick{$\ket{0}$} & \gate{H} & \gate{R_z(x_1)} & \ctrl{1} & \qw & \ctrl{1} & \gate{R_y(\theta_1)} & \ctrl{1} & \meter{} \\
\lstick{$\ket{0}$} & \gate{H} & \gate{R_z(x_2)} & \targ{} & \gate{R_z(x_1 x_2)} & \targ{} & \gate{R_y(\theta_2)} & \targ{} & \qw
\end{quantikz}
\end{center}
\end{exercise}

\begin{exercise}
\textbf{training the VQC.}
\begin{enumerate}[leftmargin=*]
  \item implement the loss function: binary cross-entropy
    \[
    \mathcal{L} = -\frac{1}{m}\sum_i \left[y_i \log(p_i) + (1-y_i)\log(1-p_i)\right]
    \]
    where $p_i = \texttt{predict}(x_i, \theta)$
  \item implement the parameter shift rule to compute gradients:
    \begin{itemize}
      \item for each parameter $\theta_k$, compute the loss at $\theta_k + \pi/2$ and $\theta_k - \pi/2$
      \item the gradient is the difference divided by 2
    \end{itemize}
  \item implement gradient descent:
    \begin{itemize}
      \item initialize $\theta$ randomly (small values, e.g., uniform in $[0, \pi]$)
      \item learning rate: start with 0.1
      \item run for 50--100 iterations
      \item record the loss at each step
    \end{itemize}
  \item plot the training loss curve. does it converge?
  \item compute test accuracy. how does it compare with the classical SVM?
\end{enumerate}

tips:
\begin{itemize}[leftmargin=*]
  \item training is slow bcs each gradient step requires $2p \times m$ circuit evaluations ($p$ parameters, $m$ training samples). use the statevector simulator to speed things up.
  \item if training is too slow, subsample the training data (use 30--40 points).
  \item try different random initializations --- variational circuits can get stuck in local minima.
\end{itemize}
\end{exercise}

\begin{exercise}
\textbf{VQC hyperparameters.}
explore how these choices affect performance:
\begin{enumerate}[leftmargin=*]
  \item number of ansatz layers (1, 2, 3). more layers $=$ more parameters $=$ more expressive, but also harder to train and more noise.
  \item feature map type: ZZ feature map vs angle encoding. which works better?
  \item number of shots (100, 1000, 10000). more shots $=$ lower variance but slower.
  \item learning rate (0.01, 0.1, 0.5). too high and it oscillates, too low and it's slow.
  \item optimizer: try both vanilla gradient descent and COBYLA (\texttt{scipy.optimize.minimize} with \texttt{method='COBYLA'}). COBYLA is gradient-free so it might work better if gradients r noisy.
\end{enumerate}
\end{exercise}

\section{the big comparison}

\begin{exercise}
\textbf{quantum vs classical showdown.}
this is the main event. compile ur results into a comparison table:

\begin{center}
\begin{tabular}{|l|c|c|}
\hline
\textbf{classifier} & \textbf{test accuracy} & \textbf{training time} \\
\hline
classical linear SVM & & \\
\hline
classical RBF SVM & & \\
\hline
QSVC (ZZ feature map, 1 rep) & & \\
\hline
QSVC (ZZ feature map, 2 reps) & & \\
\hline
VQC (1 layer) & & \\
\hline
VQC (2 layers) & & \\
\hline
\end{tabular}
\end{center}

answer these questions:
\begin{enumerate}[leftmargin=*]
  \item does the quantum classifier beat the classical one on this dataset?
  \item if not, does that mean quantum ML is useless? wt factors could change the outcome?
  \item which quantum approach (VQC or QSVC) works better? why might that be?
  \item wt's the computational cost comparison? how many circuit evaluations did QSVC require vs VQC?
  \item if u added noise (from lab 3), how would it affect each quantum method?
\end{enumerate}
\end{exercise}

\section{bonus: noisy quantum classifier}

\begin{exercise}
\textbf{running with noise.}
\begin{enumerate}[leftmargin=*]
  \item repeat the QSVC experiment but with a noisy simulator (depolarizing noise, $p = 0.01$ single-qubit, $p = 0.05$ CNOT)
  \item how much does accuracy drop?
  \item apply measurement error mitigation from lab 3 to the kernel computation
  \item does mitigation help recover the accuracy?
  \item now try ZNE on the kernel circuits. is it worth the extra computational cost?
\end{enumerate}

this exercise shows u the real challenge of near-term quantum ML: the noise can easily wipe out any advantage the quantum feature map might have.
\end{exercise}

\section{bonus: scaling up}

\begin{exercise}
\textbf{more features.}
\begin{enumerate}[leftmargin=*]
  \item repeat the QSVC experiment using all 4 iris features instead of 2 (u'll need 4 qubits)
  \item does using more features help the quantum classifier?
  \item compare the 4-feature quantum kernel with the 4-feature classical RBF kernel
  \item the 4-qubit feature map lives in a $2^4 = 16$ dimensional Hilbert space. the RBF kernel implicitly uses infinite dimensions. discuss when the quantum approach might still win despite having ``fewer'' dimensions.
\end{enumerate}

4-qubit ZZ feature map (first layer only, without interaction terms for clarity):

\begin{center}
\begin{quantikz}
\lstick{$q_0$} & \gate{H} & \gate{R_z(x_1)} & \ctrl{1} & \qw & \qw & \qw \\
\lstick{$q_1$} & \gate{H} & \gate{R_z(x_2)} & \targ{} & \ctrl{1} & \qw & \qw \\
\lstick{$q_2$} & \gate{H} & \gate{R_z(x_3)} & \qw & \targ{} & \ctrl{1} & \qw \\
\lstick{$q_3$} & \gate{H} & \gate{R_z(x_4)} & \qw & \qw & \targ{} & \qw
\end{quantikz}
\end{center}
\end{exercise}

\section{final reflection}

\begin{exercise}
\textbf{course synthesis.}
write a 1-page reflection addressing:
\begin{enumerate}[leftmargin=*]
  \item \textbf{the pipeline}: trace the full QML pipeline from data preprocessing $\to$ feature map $\to$ quantum circuit $\to$ measurement $\to$ classical post-processing. where does quantum help? where is it classical?
  \item \textbf{the bottleneck}: wt's the main practical limitation of quantum ML right now? (noise? number of qubits? training time? all of the above?)
  \item \textbf{the promise}: under wt conditions might quantum ML provide a genuine advantage? think about data structure, feature space dimensionality, and kernel expressiveness.
  \item \textbf{ur experience}: wt was the most surprising thing u learned in this course? wt would u want to explore further?
\end{enumerate}
\end{exercise}

\section{submission}

submit a jupyter notebook with:
\begin{itemize}[leftmargin=*]
  \item all exercises completed with code and output
  \item the full comparison table filled in
  \item kernel matrix heatmaps (quantum vs classical)
  \item training loss curves for the VQC
  \item decision boundary plots where possible
  \item ur final reflection
\end{itemize}

this notebook should tell a complete story: classical baseline $\to$ quantum kernel $\to$ quantum variational $\to$ comparison $\to$ reflection.

\end{document}
